\documentclass[11pt,a4paper]{article}
\usepackage[utf8]{inputenc}
\usepackage[margin=1in]{geometry}
\usepackage{hyperref}
\usepackage{enumitem}

\hypersetup{colorlinks=true, linkcolor=blue, urlcolor=blue}

\begin{document}

\noindent\textbf{To:}\\
Prof. Vincenzo Piuri, Editor-in-Chief\\
\textit{IEEE Open Journal of the Computer Society}\\
Email: \href{mailto:vincenzo.piuri@unimi.it}{vincenzo.piuri@unimi.it}\\[1.5em]

\noindent\textbf{Date:} November 20, 2026\\[1em]
\noindent\textbf{Subject:} Resubmission of Revised Manuscript (Original ID: \textbf{OJCS-2026-09-0983})\\[1em]
\noindent\textbf{Manuscript Title:}\\
\textit{Stress-Testing Preference Elicitation in Instruction-Tuned Language Models: Separating Commitment, Conditional Direction, and Payoff Sensitivity Across Small Model Families}\\
\textit{(Formerly: ``Beyond Forced Choice: Separating Commitment and Payoff Sensitivity in LLM Preference Elicitation'')}\\[1em]
\noindent\textbf{Authors:}\\
Shriyam Baloni, Achin Jain, Debojyoti Roy, Vanita Jain, and Massimo Donelli\\[1.5em]

\noindent Dear Prof. Piuri and the Editorial Board,\\[0.8em]
Thank you for the opportunity to revise and resubmit our manuscript to the \textit{IEEE Open Journal of the Computer Society}. In accordance with your decision letter dated November 6, 2026, we are resubmitting our manuscript as a new submission linked to our original manuscript ID: \textbf{OJCS-2026-09-0983}.\\[0.8em]
We have conducted a thorough, exhaustive overhaul of the paper in response to the valuable evaluations of the Associate Editor and the three reviewers. Specifically:

\begin{enumerate}[leftmargin=*]
    \item \textbf{Complete Task \& Protocol Specification (R1.1, R2.1):} Added a dedicated Methods section (\S~IV), an experimental workflow diagram (Fig.~1), formal state-transition equations, verbatim prompt phrasings (Table~II), and complete turn-by-turn transcripts (Supplementary~\S~III).
    \item \textbf{Statistical Rigor \& Separation-Robust Modeling (R2.2, R2.3, R3.2):} Explicitly formulated all joint multinomial and hurdle regression equations (Eqs.~3--7), implemented Firth penalized likelihood estimation to resolve separation on sparse cells (Table~IV, Supp.~Table~IV), aligned all point estimates with matching dilemma-cluster bootstrap intervals, and verified asymptotic convergence properties.
    \item \textbf{Threshold Follow-Up Mathematical Verification (R2.4, R3.3):} Mathematically proved that the follow-up threshold is non-binding ($s_0 + 5 \times 9 < T$ across all 20 dilemmas; \S~V-B, Supp.~\S~X) and corrected illustrative figure parameters to match the actual experimental distribution.
    \item \textbf{Matched 2AFC Baseline \& Construct Validity (R2.6, R3.1):} Evaluated a paired two-alternative forced-choice (2AFC) baseline ($N=439$), non-zero indifferent payoffs ($N=280$), and a 70-trial counterbalanced neutral-label experiment (Supp.~\S~XVI), demonstrating that simultaneously neutralizing action labels and presentation slots resolves directional disparities while preserving 100\% payoff sensitivity.
    \item \textbf{Phrasing Sensitivity \& Hardware Audit (R1.2):} Added a full cross-model phrasing breakdown (Table~VI), evaluated stochastic decoding across random seeds (Supp.~\S~XV), documented local hardware memory constraints (Supp.~\S~XI), and provided a self-contained execution suite and Jupyter notebook for evaluating unquantized FP16 and 7B models on free cloud GPUs (Supp.~\S~XVII).
    \item \textbf{Open Science \& Transparency (R1.1, R2.1, R3.6):} Fully deposited all experimental runners, deterministic response parsers, stimulus specifications, and frozen raw CSV trial records in our public GitHub repository:\\
    \url{https://github.com/CodeShrek/LLM-Preference-Diagnostics}.
\end{enumerate}

\noindent Along with this resubmission, we have provided:
\begin{itemize}[leftmargin=*]
    \item A \textbf{clean, unmarked main manuscript} (\texttt{main.pdf}), formatted strictly within the 12-page budget.
    \item A comprehensive \textbf{Supplementary Material package} (\texttt{supplementary.pdf}, 18 sections, 13 tables).
    \item A detailed, point-by-point \textbf{Response to Reviewers} document (\texttt{response\_to\_reviewers.pdf}).
\end{itemize}

\noindent We believe the manuscript is now mature, technically rigorous, and directly answers every concern raised by the reviewers. We look forward to your reconsideration.\\[1.5em]

\noindent Sincerely,\\[1em]
\textbf{Achin Jain} \textit{(Corresponding Author)}\\
Department of Computer Engineering, University of Delhi, Delhi, India\\
Email: \href{mailto:achin.mails@gmail.com}{achin.mails@gmail.com}\\[0.5em]
\textit{On behalf of co-authors:}\\
Achin Jain, Debojyoti Roy, Prof. Vanita Jain (Head, Dept. of ECE, University of Delhi), and Prof. Massimo Donelli (University of Trento)

\end{document}
